\documentclass[a4paper,10pt]{article}
\usepackage[utf8]{inputenc}
\usepackage{geometry}
\usepackage{parskip} % for spacing between paragraphs
\usepackage{hyperref} % for clickable links
\usepackage{enumitem}
\usepackage{titlesec} % to customize section titles
\usepackage{helvet} % Use Helvetica typeface
\usepackage{xcolor} % For text colors
\usepackage{ifthen} % To enable conditional statements
\usepackage{graphicx} % For the signature image
\usepackage{fancyhdr}
\usepackage{needspace} % For checking available space for sections
\usepackage[none]{hyphenat} % Disable hyphenation on word breaks

% Links: coloured text instead of hyperref's default cyan boxes
% (the boxes look dated and are meaningless in the print version).
\definecolor{linkblue}{rgb}{0.10, 0.30, 0.60}
\hypersetup{colorlinks=true, linkcolor=linkblue, urlcolor=linkblue, filecolor=linkblue}

% Section heading background for accessible layouts (overridden for inverted)
\colorlet{headingbg}{lightgray!60}

% Define a Boolean flag for accessibility features
\newboolean{accessibility}
\setboolean{accessibility}{false} % Default is off

% Handle accessibility flag from command line
\ifdefined\accessibilityflag
  \setboolean{accessibility}{true}
\fi

% Conditional font selection based on accessibility
\ifthenelse{\boolean{accessibility}}{
  % Load accessibility settings if enabled
  % Set higher contrast colors
  \definecolor{gray}{rgb}{0.2, 0.2, 0.2} % Darker gray for better contrast
  % Use a more readable line spacing
  \renewcommand{\baselinestretch}{1.5}
  
  % Define inverted color scheme
  \ifdefined\invertcolorsflag
    % Inverted high contrast colors (yellow text on black background)
    \pagecolor{black}
    \color{yellow}
    \definecolor{gray}{rgb}{0.0, 0.8, 0.8} % Cyan color for headers for better contrast with yellow text
    % Dark teal heading bar: yellow heading text on the old pale-cyan bar was unreadable
    \definecolor{headingbg}{rgb}{0.0, 0.22, 0.28}
    \hypersetup{
      colorlinks=true,
      linkcolor=cyan,
      filecolor=magenta,
      urlcolor=cyan
    }
  \fi
  
  % Check if dyslexic font flag is defined
  \ifdefined\dyslexicfontflag
    % Use OpenDyslexic font if available
    \usepackage{fontspec}
    \setmainfont[
      Path = ./fonts/,
      Extension = .otf,
      UprightFont = *-Regular,
      BoldFont = *-Bold,
      ItalicFont = *-Italic,
      BoldItalicFont = *-BoldItalic
    ]{OpenDyslexic}
  \else
    % Sans-serif (Helvetica) reads more easily, and unlike the default serif
    % font it scales, so the large name stays large.
    \renewcommand{\familydefault}{\sfdefault}
  \fi
}{
  % Standard font settings if accessibility is not enabled
  \renewcommand{\familydefault}{\sfdefault} % Set Helvetica as default font
  % Standard gray color
  \definecolor{gray}{rgb}{0.5, 0.5, 0.5}
}

% Reduce page margins
\geometry{top=0.25in, bottom=0.25in, left=0.25in, right=0.25in}

% Define a Boolean flag for dividing line (true = line, false = no line)
\newboolean{dividingline}
\setboolean{dividingline}{false} % Change to 'false' if you want to disable the line

% Define a Boolean flag for personal contact details (true = work email, no phone)
\newboolean{showworkinfo}
\setboolean{showworkinfo}{false} % Change to 'true' to show work email and hide phone

% Handle showworkinfo flag from command line
\ifdefined\showworkinfoflag
  \setboolean{showworkinfo}{true}
\fi

% Define a Boolean flag for the published link in the cover letter
\newboolean{showpublishedlink}
\setboolean{showpublishedlink}{false}

% Handle published link flag from command line
\ifdefined\showpublishedlinkflag
  \setboolean{showpublishedlink}{true}
\fi

% Custom font size for name
\newcommand{\namesection}{
    % OpenDyslexic is much wider, so a smaller size stops a long first name running into the contact block
    \ifdefined\dyslexicfontflag
        {\fontsize{50}{34}\selectfont \textbf{Taylor \newline Template}} \\%
    \else
        {\fontsize{68}{34}\selectfont \textbf{Taylor \newline Template}} \\% Font size set to 68, bold
    \fi
    \vspace{0.5em} % Small space between name and contact info
}

% Custom command for contact info with grey color and conditional phone/email display
\newcommand{\contactinfo}{
    \color{gray}\hypersetup{urlcolor=gray} Dublin, Ireland \\
    \ifthenelse{\boolean{showworkinfo}}{
        \href{mailto:jobs@example.com}{jobs@example.com} \\
    }{
        \href{tel:+353000000000}{+353 00 000 0000} \\
        \href{mailto:taylor@example.com}{taylor@example.com} \\
    }
    \href{https://example.com}{example.com} \\
    \href{https://github.com/benkelly/cv-pipeline}{github.com/benkelly/cv-pipeline} \\
}

% ponytail: the name block is body content on page 1, not a running header. As a
% fancyhdr running header it repeated on every page, and its 133.92pt headheight
% reserved that much empty top margin on every page whether it drew or not.
\pagestyle{empty}
\newcommand{\namebanner}{%
    \begin{minipage}[t]{0.5\textwidth}%
        \namesection
    \end{minipage}%
    \begin{minipage}[t]{0.5\textwidth}%
        \raggedleft
        \contactinfo
    \end{minipage}%
    \par\vspace{1em}
}

% Use the plain style for the cover letter
% \thispagestyle{plain}


% Signature image. Optional: drop your own static/sig.png in to use it.
\newcommand{\sig}{
    \IfFileExists{static/sig.png}{
        \includegraphics[width=0.18\textwidth]{static/sig.png}
        \ifthenelse{\boolean{accessibility}}{
            % Add accessibility note for screen readers
            \mbox{}\\
            \textit{\small [Image: handwritten signature]}
        }{}
    }{}
}

% Optional published link for cover letter builds
\providecommand{\PUBLISHEDURL}{}
\newcommand{\publishedlink}{
    \ifthenelse{\boolean{showpublishedlink}}{
        \leavevmode\newline
        \vspace{1em}
        {\small\color{gray}Full CV and cover letter available at \href{\PUBLISHEDURL}{\PUBLISHEDURL}}
    }{}
}

% Define a reusable CV section command with three arguments, optional fourth argument, and optional dividing line
\newcommand{\cvsection}[4][]{
    \ifthenelse{\boolean{accessibility}}{
        \ifdefined\newpagesectionsflag
            % Start each section on a new page with no header
            % \newpage
            % \thispagestyle{empty}
            \newpage \thispagestyle{plain} \setlength{\headheight}{0pt} 
            % Check if there's enough space (approximately 5cm) for the section
            % If not enough space, a page break will be inserted automatically
        % \else
        %     % Check if there's enough space (approximately 5cm) for the section
        %     \needspace{5cm}
        %     % If not enough space, a page break will be inserted automatically
        \fi
        % Only draw the rule and heading bar for real section titles; untitled
        % sub-entries (e.g. later roles at the same company) used to get an empty grey bar.
        \ifthenelse{\equal{#2}{}}{
            \vspace{1em}
        }{
            % Keep the heading with its first lines instead of stranding it at a page bottom
            \needspace{4\baselineskip}
            \vspace{1.5em} % More space between sections for accessibility
            \noindent\rule{\textwidth}{2pt} % Thicker dividing line for accessibility
            \vspace{0.5em} % Space after the dividing line
            % Add section title in a distinctive background box for better visibility
            \noindent\colorbox{headingbg}{
                \begin{minipage}{\dimexpr\textwidth-2\fboxsep\relax}
                    \vspace{0.2em}
                    \raggedright\textbf{\Large #2}
                    \vspace{0.2em}
                \end{minipage}
            }
            \vspace{0.5em}
        }
    }{
        % Check if there's enough space (approximately 5cm) for the section
        \needspace{5cm}
        % If not enough space, a page break will be inserted automatically
        \vspace{1em} % Standard spacing between sections
        \ifthenelse{\boolean{dividingline}}{
            \noindent\rule{\textwidth}{0.5pt} % If true, add a dividing line
        }{}
        \begin{minipage}[t]{0.25\textwidth} % Section title on the left (25%)
            \raggedright \color{gray}\textbf{#2}
        \end{minipage}%
        \begin{minipage}[t]{0.75\textwidth} % Section content on the right (75%)
            \raggedright \textbf{#3} \\
            #4
            \ifthenelse{\equal{#1}{}}{}{\\{\small\color{gray}#1}} % Optional argument displayed if provided
        \end{minipage}
    }
    \ifthenelse{\boolean{accessibility}}{
        % For accessible versions, use a different layout with section title at top
        \begin{minipage}[t]{\textwidth} % Content below the section title
            \raggedright \textbf{#3} \\
            #4
            \ifthenelse{\equal{#1}{}}{}{\\{\small\color{gray}#1}} % Optional argument displayed if provided
        \end{minipage}
    }{}
    \vspace{1em}
}

% Custom command to render cover letter (with no fancy header)
\newcommand{\coverletter}[2]{
    % Name and Contact Info on the same line
    % \noindent
    \leavevmode\newline
    % \leavevmode\newline
    \vspace{1em} % Some space before the cover letter content
    #1 % Recipient information (e.g., company, address, date)
    \leavevmode\newline
    \vspace{1em} % Space before the body of the letter
    #2 % Body of the cover letter
    \vspace{1em} % Space before the signature
    % \noindent
    % \leavevmode\newline
    % \leavevmode\newline
    % \noindent
    % \leavevmode\newline
    \sig % Insert the signature (SVG)
    \leavevmode\newline
    Taylor Template

    
    \publishedlink

    \leavevmode\newpage % Insert a page break after the cover letter
}

% % Allow passing metadata from command line
% \usepackage{shellesc} % To enable shell escape for git commands
% \providecommand{\VERSION}{\detokenize{\ShellEscape{git describe --tags --long 2>/dev/null || echo "v0.0.1-$(git rev-parse --short HEAD)"}}}

% Allow passing metadata from command line
% Define VERSION command (default if not provided)
\providecommand{\VERSION}{local-build}

% Simple date format instead of using datetime package
\providecommand{\simpledate}{\number\year-\ifnum\month<10 0\fi\number\month-\ifnum\day<10 0\fi\number\day}

% PDF metadata - these are provided as defaults that can be overridden by \def commands from pdflatex
\providecommand{\PDFTITLE}{Taylor Template CV -- \simpledate{} -- \VERSION}
\providecommand{\PDFAUTHOR}{Taylor Template}
\providecommand{\PDFSUBJECT}{Taylor Template CV -- \simpledate{} -- \VERSION}
\providecommand{\PDFKEYWORDS}{CV, Taylor Template, Platform Engineer, Kubernetes, Terraform, Go, Python}

% Set PDF metadata using hyperref
\hypersetup{
    pdftitle={\PDFTITLE},
    pdfauthor={\PDFAUTHOR},
    pdfsubject={\PDFSUBJECT},
    pdfkeywords={\PDFKEYWORDS, \VERSION},
    pdfcreator={cv-pipeline, github.com/benkelly/cv-pipeline},
    pdfproducer={LaTeX}
}

% Define content input
\begin{document}

% Determine what content to include
\namebanner

\ifdefined\coverletteronly
    % Just include the cover letter with the name section
    % coverletter.tex - optional cover letter.
%
% To turn the cover letter off, comment out the \coverletter{...}{...} block.
% The pipeline notices and builds the CV on its own.

\def\recipient{%
    \textbf{Hiring Team} \\
    Example Corp, Dublin \\
    \textbf{Date:} \today \\
}

\coverletter{\recipient}{
Dear Hiring Team, \\

I'm applying for the Platform Engineer role at Example Corp. I build platforms that let product teams ship safely without waiting on anyone, and I'd love to do that for your teams too. \\

At Example Cloud I lead the team behind our internal developer platform, used by 40 teams every day. Before that I ran payment services across three regions at Acme Payments. \\

Thank you for your time and consideration. \\

Best regards, \\
}

\else
    \ifdefined\cvonly
        % Just include the CV content without cover letter
        % cv.tex - all CV content goes here.
%
% Everything in this file is made up. Replace it with your own.
%
% \cvsection[note]{Section title}{Heading}{Body}
%   - Section title: shown in the left column (leave empty for a follow-on entry)
%   - Heading: bold line, e.g. job title and company
%   - Body: text or an itemize list
%   - note (optional): small grey line under the entry

%%%%%%%%%%%%%%%%%%%
% Summary Section %
%%%%%%%%%%%%%%%%%%%
\cvsection{Summary}{}{
    Platform engineer with seven years' experience building cloud infrastructure, developer platforms and reliable services. Enjoys making the safe path the easy path for product teams.
}

%%%%%%%%%%%%%%%%%%%%%%
% Experience Section %
%%%%%%%%%%%%%%%%%%%%%%
\cvsection{Experience}{Senior Platform Engineer, Example Cloud Ltd}{
    Dublin, Ireland \\
    March 2023 -- Present \\
    \begin{itemize}
        \item Lead a team of four running the internal developer platform on Kubernetes for 40 product teams.
        \item Moved infrastructure to reusable Terraform modules, cutting new-service setup from two weeks to one day.
        \item Introduced service-level objectives and on-call runbooks, reducing paging alerts by 60\%.
        \item Built a self-service preview environment for every pull request.
    \end{itemize}
}

\cvsection{}{Software Engineer, Acme Payments}{
    Cork, Ireland \\
    June 2019 -- February 2023 \\
    \begin{itemize}
        \item Built and operated Go services handling card authorisations across three regions.
        \item Designed the multi-region failover plan and ran quarterly game days to test it.
        \item Wrote observability standards with Prometheus, Grafana and OpenTelemetry.
    \end{itemize}
}

\cvsection{}{Graduate Engineer, Sample Analytics}{
    Galway, Ireland \\
    September 2018 -- May 2019 \\
    \begin{itemize}
        \item Maintained batch data pipelines in Python and Apache Spark.
    \end{itemize}
}

%%%%%%%%%%%%%%%%%%
% Skills Section %
%%%%%%%%%%%%%%%%%%
\cvsection{Skills}{}{
    \textbf{\color{gray}Languages:} Go, Python, TypeScript, Bash, SQL \\
    \textbf{\color{gray}Cloud \& Infrastructure:} AWS, GCP, Kubernetes, Terraform, Helm, Docker, GitOps \\
    \textbf{\color{gray}Reliability:} Prometheus, Grafana, OpenTelemetry, Incident Response, Post-mortems \\
    \textbf{\color{gray}Security:} RBAC, Just-in-Time Access, PKI, Secrets Management
}

%%%%%%%%%%%%%%%%%%%%
% Projects Section %
%%%%%%%%%%%%%%%%%%%%
\cvsection{Projects}{}{
    \textbf{\href{https://github.com/benkelly/cv-pipeline}{cv-pipeline}}: This CV. One LaTeX file builds standard, accessible, dyslexia-friendly and high-contrast PDFs, plus Markdown, JSON and HTML.
}

%%%%%%%%%%%%%%%%%%%%%
% Education Section %
%%%%%%%%%%%%%%%%%%%%%
\cvsection{Education}{Example University - BSc Computer Science, 2018}{
    \textit{Thesis}: ``Scheduling batch jobs on spare cloud capacity''
}

    \else
        % Include both - first cover letter, then CV content
        % coverletter.tex - optional cover letter.
%
% To turn the cover letter off, comment out the \coverletter{...}{...} block.
% The pipeline notices and builds the CV on its own.

\def\recipient{%
    \textbf{Hiring Team} \\
    Example Corp, Dublin \\
    \textbf{Date:} \today \\
}

\coverletter{\recipient}{
Dear Hiring Team, \\

I'm applying for the Platform Engineer role at Example Corp. I build platforms that let product teams ship safely without waiting on anyone, and I'd love to do that for your teams too. \\

At Example Cloud I lead the team behind our internal developer platform, used by 40 teams every day. Before that I ran payment services across three regions at Acme Payments. \\

Thank you for your time and consideration. \\

Best regards, \\
}

        % cv.tex - all CV content goes here.
%
% Everything in this file is made up. Replace it with your own.
%
% \cvsection[note]{Section title}{Heading}{Body}
%   - Section title: shown in the left column (leave empty for a follow-on entry)
%   - Heading: bold line, e.g. job title and company
%   - Body: text or an itemize list
%   - note (optional): small grey line under the entry

%%%%%%%%%%%%%%%%%%%
% Summary Section %
%%%%%%%%%%%%%%%%%%%
\cvsection{Summary}{}{
    Platform engineer with seven years' experience building cloud infrastructure, developer platforms and reliable services. Enjoys making the safe path the easy path for product teams.
}

%%%%%%%%%%%%%%%%%%%%%%
% Experience Section %
%%%%%%%%%%%%%%%%%%%%%%
\cvsection{Experience}{Senior Platform Engineer, Example Cloud Ltd}{
    Dublin, Ireland \\
    March 2023 -- Present \\
    \begin{itemize}
        \item Lead a team of four running the internal developer platform on Kubernetes for 40 product teams.
        \item Moved infrastructure to reusable Terraform modules, cutting new-service setup from two weeks to one day.
        \item Introduced service-level objectives and on-call runbooks, reducing paging alerts by 60\%.
        \item Built a self-service preview environment for every pull request.
    \end{itemize}
}

\cvsection{}{Software Engineer, Acme Payments}{
    Cork, Ireland \\
    June 2019 -- February 2023 \\
    \begin{itemize}
        \item Built and operated Go services handling card authorisations across three regions.
        \item Designed the multi-region failover plan and ran quarterly game days to test it.
        \item Wrote observability standards with Prometheus, Grafana and OpenTelemetry.
    \end{itemize}
}

\cvsection{}{Graduate Engineer, Sample Analytics}{
    Galway, Ireland \\
    September 2018 -- May 2019 \\
    \begin{itemize}
        \item Maintained batch data pipelines in Python and Apache Spark.
    \end{itemize}
}

%%%%%%%%%%%%%%%%%%
% Skills Section %
%%%%%%%%%%%%%%%%%%
\cvsection{Skills}{}{
    \textbf{\color{gray}Languages:} Go, Python, TypeScript, Bash, SQL \\
    \textbf{\color{gray}Cloud \& Infrastructure:} AWS, GCP, Kubernetes, Terraform, Helm, Docker, GitOps \\
    \textbf{\color{gray}Reliability:} Prometheus, Grafana, OpenTelemetry, Incident Response, Post-mortems \\
    \textbf{\color{gray}Security:} RBAC, Just-in-Time Access, PKI, Secrets Management
}

%%%%%%%%%%%%%%%%%%%%
% Projects Section %
%%%%%%%%%%%%%%%%%%%%
\cvsection{Projects}{}{
    \textbf{\href{https://github.com/benkelly/cv-pipeline}{cv-pipeline}}: This CV. One LaTeX file builds standard, accessible, dyslexia-friendly and high-contrast PDFs, plus Markdown, JSON and HTML.
}

%%%%%%%%%%%%%%%%%%%%%
% Education Section %
%%%%%%%%%%%%%%%%%%%%%
\cvsection{Education}{Example University - BSc Computer Science, 2018}{
    \textit{Thesis}: ``Scheduling batch jobs on spare cloud capacity''
}

    \fi
\fi

\end{document}
